\conclusion

В рамках данной работы была спроектирована и реализована распределённая система
вычислений, в основе которой лежит кластер узлов,
использующих алгоритм консенсуса Raft для поддержания согласованного состояния.
Предложенная архитектура разделяет систему на контур координации
и контур исполнения: кластер Raft отвечает за приём клиентских запросов,
репликацию и фиксацию изменений состояния, а вычислительные узлы выполняют
пользовательские задачи и возвращают результаты через согласованный
серверный контур. Такое разделение позволило совместить требования
к отказоустойчивости, детерминированности и масштабируемости.

В работе была детально описана архитектура основных компонентов системы:
Raft-узлов, клиентской части и вычислительных узлов. Для реализации сетевого
взаимодействия был использован стек \texttt{gRPC + Protocol Buffers},
обеспечивающий строгую типизацию интерфейсов и удобство сопровождения.
Отдельное внимание было уделено реализации журнала Raft, машины состояний,
механизма снимков данных, серверных сервисов и подсистемы диспетчеризации задач.
В результате была получена рабочая программная система, поддерживающая
распределённую постановку и выполнение вычислительных задач.

Существенной частью работы стало формальное описание алгоритма консенсуса
на языке TLA+. Разработанная спецификация охватывает основные элементы Raft:
состояния узлов, термы, голосование, репликацию журнала, продвижение
\texttt{commitIndex}, обработку RPC-сообщений и ключевые инварианты
безопасности. Формальная модель позволила зафиксировать не только ожидаемое
поведение системы, но и те свойства, нарушение которых недопустимо
в корректной реализации.

Наиболее важным результатом работы стала разработка и практическое применение
подхода к проверке соответствия реализации и спецификации.
Для этого реализация была инструментирована встроенным трассировщиком,
фиксирующим абстрактные события исполнения в формате NDJSON.
Далее был построен автоматизированный конвейер проверки, включающий запуск
кластера, сбор и объединение трасс, а также их последующую валидацию
модельным проверяющим TLC через обёртку \texttt{RaftTrace.tla}.
Тем самым формальная спецификация использовалась не только как средство
описания алгоритма, но и как практический критерий корректности
реальной реализации.

Проведённая проверка показала, что связка
``инструментированная реализация --- трасса исполнения --- TLA+-проверка''
является эффективным способом валидации распределённого кода.
В отличие от обычных модульных и интеграционных тестов, такой подход
сопоставляет наблюдаемое поведение программы с полным множеством переходов,
разрешённых спецификацией, и позволяет выявлять ошибки,
проявляющиеся только при определённых межпоточных и сетевых сценариях.
С помощью разработанного инструментария были обнаружены и локализованы
два подтверждённых расхождения реализации со спецификацией:
неверное начальное значение терма и наличие self-peer в рассылке
\texttt{RequestVote}. Оба дефекта затрагивали базовую логику алгоритма Raft,
что подтверждает практическую полезность выбранного метода проверки.

Таким образом, поставленные в работе задачи были выполнены:
разработана архитектура распределённой системы вычислений,
реализован её программный прототип, построена формальная спецификация
ключевого алгоритма консенсуса и создана инфраструктура для проверки
соответствия реализации этой спецификации.
Полученные результаты показывают, что сочетание инженерной реализации
и формальных методов позволяет не только повысить доверие к корректности
системы, но и использовать спецификацию как рабочий инструмент анализа
и отладки.

Перспективы дальнейшего развития работы связаны прежде всего с расширением
разработанной библиотеки трассировки и переносом предложенного подхода на более
широкий круг систем. В ближайших планах находится поддержка новых языков
программирования в библиотеке трассировки; минимальной целевой точкой является
добавление языка Lua, что особенно важно для встраивания трассировки в СУБД
Tarantool. Такая интеграция позволит применять подход ``инструментированная
реализация --- трасса исполнения --- формальная проверка'' не только к учебным
или исследовательским реализациям распределённых алгоритмов, но и к
промышленным системам хранения и обработки данных.
